% Immutable cited-input supplement. Target proofs remain embargoed.
% Target: wbII:prop:leray-edge-restriction; tag 03IM.
% Frozen source SHA256 74497bfe327fae25c9cdcb17334b3bf42a9793a9df7c0a4404d97568f2f760f0.
% CompendiumRef is resolved against the cited S6_09 source edition, not silently identified with the current series statement.

% Exact cited input: comp:I / mapsI:top:prop:leray-e2
\begin{proposition}[all Leray terms used below]
\label{mapsI:top:prop:leray-e2}
\ The specialization theorem
\cref{mapsI:top:thm:p0-integral-specialization}, the finite fillings of
\cref{wb:prop:candidate-finite-fillings}, and
\cref{wb:lem:finite-filling-stalk-map} give, for \(j=1,2\) and
\(0\le r\le4\), the finite-filling compatibility
\[
\tag{IL.FF}
 N_j\ \text{strongly deformation retracts onto }S_j,\qquad
 H^r(S_j;\mathbb Z)\longrightarrow H^r(F;\mathbb Z)^{T_j}
 \ \text{is }\pi_j^*
\]
under the marked punctured-fibre identification.  Consequently,
\[
H^0(B,R^1)=\mathbb Z(12\gamma),\quad
H^0(B,R^2)=\mathbb Z(2q),\quad
H^0(B,R^3)=\mathbb Z(2\gamma uw),
\tag{IL.5.10}
\]
\[
 H^1(B,R^1)=H^1(B,R^2)=0,
\tag{IL.5.11}
\]
and, with \(\omega\) a generator of \(H^2(B;\mathbb Z)\),
\[
H^2(B,R^1)=\mathbb Z\omega[\delta],\qquad
H^2(B,R^2)=\mathbb Z\omega[\gamma\delta],\qquad
H^2(B,R^4)=\mathbb Z\omega[\operatorname{vol}].
\tag{IL.5.12}
\]
These identifications are compatible with exterior products in \(V\).
\end{proposition}

% Exact cited input: comp:I / mapsI:top:lem:higher-transgressions
\begin{lemma}[the higher transgressions]
\label{mapsI:top:lem:higher-transgressions}
\ In the integral Leray spectral sequence of the constructed map,
\[
 d_2^{0,2}(2q)=\omega[\delta],
 \qquad
 d_2^{0,3}(2\gamma uw)=\omega[\gamma\delta].
\tag{IL.5.22}
\]
\end{lemma}

% Exact cited input: comp:I / mapsI:top:thm:p0-integral-specialization
\begin{theorem}[nearby-cycle specialization]
\label{mapsI:top:thm:p0-integral-specialization}
\ For the proper holomorphic cusp filling
\(f_0:N_0\to\Delta\) of \cref{wb:thm:candidate-cusp-geometry}, with its
proved local normal-crossings equations and normalization combinatorics,
the specialization map of the central fibre \(W\),
\[
 \operatorname{sp}^r:H^r(W;\mathbb Z)\longrightarrow H^r(F;\mathbb Z)
\]
is injective and
\[
 \operatorname{im}\operatorname{sp}^r
 =H^r(F;\mathbb Z)^{T_0}
\tag{IL.5.1}
\]
for \(0\le r\le4\).  In particular, the image is saturated and there is no
unrecorded finite-index defect at \(p_0\).
\end{theorem}

% Exact cited input: comp:I / wb:prop:candidate-finite-fillings
\begin{proposition}[finite fillings]
\label{wb:prop:candidate-finite-fillings}
\ Let \(A_j\) be the two finite monodromies and choose
\[
 v_1=\varepsilon,\qquad v_2=-\varepsilon',
 \qquad \gamma(v_1)=1,\quad\gamma(v_2)=-1.
\]
On the local torus family over a linearizing disc \(\Delta_j\), the affine
map
\[
 (z,x)\longmapsto(g_jz,A_jx+v_j/m_j)
\tag{2.9}
\]
generates a free action of order \(m_j\).  The quotient
\(f_j:N_j\to D_j\) is a smooth complex threefold with
\[
 f_j^*(p_j)=m_jS_j,\qquad
 S_j=A_0^{(j)}/\langle x\mapsto A_jx+v_j/m_j\rangle.
\tag{2.10}
\]
Over the punctured disc it is biholomorphic to the original period family.
\end{proposition}

% Exact cited input: comp:I / wb:lem:finite-filling-stalk-map
\begin{lemma}[radial retraction and the finite-fibre stalk map]
\label{wb:lem:finite-filling-stalk-map}
\ Under \cref{wb:prop:candidate-finite-fillings}, the finite local
model \(N_j\) strongly deformation retracts onto its reduced central fibre
\(S_j\).  After identifying a nearby fibre with the marked real torus
\(A_s=(\Lambda\otimes\mathbb R)/\Lambda\) by the flat coordinates
\(\Xi\), the specialization map
\[
 H^r(S_j;\mathbb Z)\longrightarrow H^r(A_s;\mathbb Z)
\]
is the pullback \(\pi_j^*\) of the quotient covering
\[
 \pi_j:A_0^{(j)}\longrightarrow
 S_j=A_0^{(j)}/\langle x\mapsto A_jx+v_j/m_j\rangle .
\]
\end{lemma}

% Exact cited input: comp:I / wb:thm:candidate-cusp-geometry
\begin{theorem}[the cusp quotient and its central fibre]
\label{wb:thm:candidate-cusp-geometry}
\ For the period functions and holomorphic cusp matrix \(C(t)\)
constructed above, and sufficiently small \(\epsilon>0\), the maps
\[
 \Psi_{\bar\lambda}(x,t)=
 \bigl(c_{\bar\lambda}(t)t^{B_0\bar\lambda}x,t\bigr),
 \qquad
 c_{\bar\lambda}(t)=\exp(2\pi iC(t)\bar\lambda),
\tag{2.5}
\]
extend to a free properly discontinuous action of
\(\bar\Lambda\simeq\mathbb Z^2\) on
\(Y_\epsilon=\{|t|<\epsilon\}\).  The quotient
\[
 f_0:N_0=Y_\epsilon/\bar\Lambda\longrightarrow\Delta_\epsilon
\]
is a proper holomorphic map from a connected smooth complex threefold.
Its nonzero fibres are the tori of \cref{wb:prop:period-input-audit}.  Its full
scheme-theoretic central fibre \(W\) is reduced and has local equations
\[
 z_0,\qquad z_0z_1,\qquad z_0z_1z_2.
\tag{2.6}
\]
The surface \(W\) is irreducible and nonnormal.  Its normalization is
\(\nu:E_0\simeq\mathrm{dP}_6\to W\), obtained by identifying opposite sides
of the toric hexagon; its double locus consists of three rational curves
through exactly two triple points.  Thus \(W\) is a normal-crossings surface,
but it is not a simple-normal-crossings divisor.  Moreover
\[
 \omega_{N_0}\simeq\mathcal O_{N_0}.
\tag{2.6a}
\]
\end{theorem}

% Exact cited input: comp:I / wb:prop:period-input-audit
\begin{proposition}[existence and non-degeneracy of the period matrix]
\label{wb:prop:period-input-audit}
\ There are holomorphic functions
\[
 \tau:\mathfrak h_z\to\mathfrak h,
 \qquad \mu,\beta:\mathfrak h_z\to\mathbb C
\]
satisfying
\[
\begin{array}{lll}
\tau\circ g_0=\tau-1,&\mu\circ g_0=\mu,&\beta\circ g_0=\beta+1,\\
\tau\circ g_2=-1/\tau,&\mu\circ g_2=1+\mu/\tau,
 &\beta\circ g_2=\beta-3-6\mu^2/\tau,
\end{array}
\tag{2.1}
\]
and the corresponding order-three transformation laws obtained from
\(\tau\circ g_1=(\tau-1)/\tau\).  They may be chosen so that
\[
 \operatorname{Im}\tau>0,
 \qquad
 D:=\operatorname{Im}\beta-
 \frac{6(\operatorname{Im}\mu)^2}{\operatorname{Im}\tau}<0.
\tag{2.2}
\]
Consequently
\[
 \Pi=[Z\mid I_2],\qquad
 Z=\begin{pmatrix}6\mu&\tau\\ \beta&\mu\end{pmatrix}
\tag{2.3}
\]
maps the marked rank-four lattice \(\Lambda\) isomorphically onto a lattice
in \(\mathbb C^2\) at every point.  The quotients
\(\mathbb C^2/\Pi(z)\Lambda\) form a holomorphic family of compact
two-dimensional complex tori over \(\mathfrak h_z\), equivariant for
\(\Delta\).
\end{proposition}

